% Discussion paragraph for Storyline 1: the two routes out of the punishment dilemma.
% In the manuscript, replace Eq.~(3) by Eq.~\eqref{eq:lambda-slope} and Fig.~5a by Fig.~\ref{fig:opportunism}a.
Our analysis suggests a general way to organize remedies for the punishment dilemma. The harm that disputed judgements do to justified punishers depends on how often such judgements occur and on how strongly each one changes a reputation. Institutions, gossip, synchronization and generous assessment act on the first factor\cite{radzvilavicius2021,kawakatsu2024,murase2024synchronization,schmid2021}, whereas scores, aggregation and inertia act on the second\cite{schmid2023,michelmata2024}. The same separation appears in our theory. Near a resident population, the reputational cost of withholding help is the product of a term from the first step, the assessment contrast, which measures how differently observers judge cooperation and defection, and a term from the second step, the slope of the reputation map, which measures how strongly a change in assessments moves a reputation (Eq.~(3)). The two routes are therefore complementary. Weak agreement can be offset by robust revision, as in our model, and fragile revision does little harm when information is public. Leniency, however, has a price at either step. Generous assessment and tolerant revision both protect deliberate defectors as well as justified punishers\cite{schmid2021}, as our comparison of inertia, aggregation and scores shows (Fig.~5a). A reputation system should therefore separate whether a single negative assessment can overturn a reputation, which inertia suppresses, from what it costs at the margin, which partial consequences restore. Disputed evidence should be costly but not catastrophic.
